% Matrix-chain multiplication: the result is the same whatever order you
% choose, but the cost is not -- which is what dynamic programming exploits.
\documentclass[dvisvgm,border=3pt]{standalone}
\input{preamble.tex}
\begin{document}
\begin{tikzpicture}[x=1cm,y=1cm,
    lbl/.style={text=fg,font=\sffamily\small,align=center}]

  \def\u{0.135}   % one matrix row/column

  % \mat{x}{rows}{cols}{name}{dims}
  % The matrices differ in height, so the captions sit on two fixed baselines
  % below the tallest one (8 rows) rather than below each matrix.
  \def\capa{-1.42}
  \def\capb{-1.76}
  \newcommand{\mat}[5]{%
    \pgfmathsetmacro{\w}{#3*\u}
    \foreach \c in {0,...,\the\numexpr#3-1\relax}{
      \foreach \r in {0,...,\the\numexpr#2-1\relax}{
        \draw[draw=uibkorange!70!black,line width=0.35pt,fill=uibkorange]
          ({#1+\c*\u},{-\r*\u}) rectangle ++(\u,-\u);}}
    \node[lbl] at ({#1+\w/2},\capa) {#4};
    \node[lbl,font=\sffamily\scriptsize] at ({#1+\w/2},\capb) {#5};}

  \mat{0.00}{5}{5}{A}{5x5}
  \node[lbl] at (0.82,-0.34) {$\times$};
  \mat{1.00}{5}{4}{B}{5x4}
  \node[lbl] at (1.68,-0.34) {$\times$};
  \mat{1.86}{4}{8}{C}{4x8}
  \node[lbl] at (3.10,-0.34) {$\times$};
  \mat{3.28}{8}{2}{D}{8x2}
  \node[lbl] at (3.72,-0.34) {$=$};
  \mat{3.90}{5}{2}{X}{5x2}

  % the three ways to parenthesise, and what each costs
  \node[lbl,anchor=west,font=\sffamily\small] at (0,-2.35)
    {$A \times \bigl(B \times (C \times D)\bigr) = X$};
  \node[lbl,anchor=west,font=\sffamily\small] at (0,-2.92)
    {$(A \times B) \times (C \times D) = X$};
  \node[lbl,anchor=west,font=\sffamily\small] at (0,-3.49)
    {$\bigl((A \times B) \times C\bigr) \times D = X$};
  \node[lbl,anchor=west] at (3.95,-2.35) {154 ops};
  \node[lbl,anchor=west] at (3.95,-2.92) {204 ops};
  \node[lbl,anchor=west] at (3.95,-3.49) {340 ops};

\end{tikzpicture}
\end{document}
